
We investigated whether execution-based code generation benchmarks (HumanEval, MBPP) produce distinctive evaluation signatures. Data infrastructure was successfully established with 100\% feature completeness across 14 model-benchmark pairs. Analysis revealed perfect ranking correlation (Spearman ρ = 1.000, p < 0.0001) between benchmarks despite substantial distributional divergence (KL = 18.4).

This finding indicates that HumanEval and MBPP, while differing in statistical properties (difficulty calibration), produce identical model competency orderings. The result does not support the hypothesis that benchmark design differences create distinctive evaluation signatures. The most plausible interpretation is that execution-based benchmarks measure a unidimensional competency construct (general code generation ability) with different difficulty scaling.

Limitations include small sample size (n=6 models), restriction to two benchmarks, and incomplete analysis chain (2 of 5 planned sub-analyses completed). Results should be validated with larger model populations and broader benchmark coverage.

The finding informs benchmark selection: using multiple execution-based benchmarks may provide redundant ranking information. Evaluation diversity likely requires task-type variety (generation, understanding, repair, translation) rather than multiple instances of the same task type. Future benchmark design should prioritize cross-task-type coverage over within-task-type redundancy.
